% ============================================================
% CHAPTER 6 — Results (REVISED)
% ============================================================

\chapter{Results}
\label{chap:results}

This chapter presents the experimental results obtained from the two
benchmark evaluations described in Chapter~\ref{chap:experimental_setup}.
The results are organised to address the four evaluation questions
defined in Section~\ref{sec:experimental_objectives}.
Section~\ref{sec:expense_results} reports on the Expense Report benchmark,
covering Pareto recovery correctness and compliance behaviour
(\textbf{EQ1} and \textbf{EQ2}).
Section~\ref{sec:results_loanapp} reports on the role-conditioned Loan
Application benchmark, evaluating scalability and comparative positioning
against prior work (\textbf{EQ3} and \textbf{EQ4}).

% ============================================================
\section{Exact Pareto Recovery: Expense Report}
\label{sec:expense_results}

This section addresses \textbf{EQ1} (correctness) and
\textbf{EQ2} (compliance guarantee). The Expense Report
benchmark provides a controlled setting in which the exact
Pareto front is known from an independent CSP+COP solver,
enabling a direct correctness comparison between the proposed
\ac{RL} framework and an exact optimisation method.

\subsection{Ground Truth and Baseline Comparison}
\label{sec:expense_baseline}

As ground truth, the CSP+COP solver was executed on the Expense Report
\ac{DCR} graph with maximum trace length $k=5$, minimising both cost and duration simultaneously.
The solver identifies two Pareto-optimal traces in 1.01 seconds
(Table~\ref{tab:expense_pareto_csp}).

Following Diaz et al. \cite{Diaz2024Pareto-OptimalModels}, cost values are expressed in euros (€) and duration in minutes (min); these are the only experiments in this thesis where concrete units are assigned to the objective dimensions.

\begin{table}[H]
\centering
\caption{Exact Pareto front obtained with
         CSP+COP~\cite{Diaz2024Pareto-OptimalModels}.}
\label{tab:expense_pareto_csp}
{\renewcommand{\arraystretch}{1.5}
\setlength{\tabcolsep}{8pt}
\begin{tabular}{lrrl}
\toprule
\textbf{Point} & \textbf{Cost (€)} & \textbf{Duration (min)} & \textbf{Trace} \\
\midrule
P1 & 50  & 10{,}500 & FillOut $\rightarrow$ Approve $\rightarrow$ PayOut \\
P2 & 710 & 2{,}700  & FillOut $\rightarrow$ Withdraw \\
\bottomrule
\end{tabular}}
\end{table}

A PPO agent was trained using a normalised reward formulation
(Section~\ref{sec:reward_normalised}).
Under this formulation the agent recovers both Pareto-optimal points, matching the exact solver output.

Reward normalisation (Section~\ref{sec:reward_normalised}) is essential
for stable learning in this setting: without it, a single high-cost
action's raw penalty can dominate the episodic signal and overwhelm
the gradient that guides the agent toward accepting states.

The reward convergence curve under the normalised formulation is shown in Figure~\ref{fig:reward_convergence}.

\begin{figure}[H]
    \centering
    \includegraphics[width=0.8\linewidth]{Pictures/Expense_Report_reward_convergence.png}
    \caption{Reward convergence across weight configurations after
             reward normalisation.}
    \label{fig:reward_convergence}
\end{figure}

Across all six configurations, the mean per-episode reward rises
sharply during the first decile of training (from approximately $89$
to $98$) and then stays within a narrow band ($96$--$101$) for the
remaining $90\,\%$ of training (per-decile means in
Appendix~\ref{app:expense_reward_deciles}). The final $20\,\%$ window
used throughout this thesis for Pareto and compliance statistics
(Section~\ref{sec:training_pipeline}) therefore lies well inside this
converged regime, with substantial margin.

% ------------------------------------------------------------
\subsection{Compliance and Illegal-Action Ratio}
\label{sec:expense_compliance}

The shielding mechanism operates at the environment level: any action not enabled under the current \ac{DCR} marking is intercepted before execution, leaving the process state unchanged and returning a fixed
penalty. This guarantee holds throughout training, the process never enters an invalid state regardless of agent behaviour.
The illegal-action ratio therefore measures how frequently the agent \emph{attempted} blocked actions, and reflects learning dynamics rather than any compliance violation.

Across the full training run, between 4.4\% and 8.7\% of all steps
are illegal action attempts, a ratio that reflects early exploration
of the action space rather than systematic non-compliance. A
complementary, episode-level view of the converged phase tells the
same story: with the normalised reward, fewer than 15\,\% of episodes
in the final 20\,\% of training contain even a single illegal
attempt (Table~\ref{tab:illegal_actions}), and the acceptance rate
reaches 100\,\% in all cases.


\begin{table}[H]
\centering
\caption{Distribution of episodes by number of illegal-action attempts within them, final 20\,\% of training, normalised
         reward. Each row is one weight pair;}
\label{tab:illegal_actions}
{\renewcommand{\arraystretch}{1.3}
\setlength{\tabcolsep}{6pt}
\begin{tabular}{lrrr}
\toprule
\textbf{Weight pair} & \shortstack{\textbf{0 illegal}\\\textbf{(\% episodes)}} & \shortstack{\textbf{1 illegal}\\\textbf{(\% episodes)}}
& \shortstack{\textbf{2+ illegal}\\\textbf{(\% episodes)}} \\
\midrule
$\alpha=0.0,\,\beta=0.0$ & 99.5 & 0.5 & 0.0 \\
$\alpha=1.0,\,\beta=0.0$ & 85.9 & 7.5 & 6.7 \\
$\alpha=0.0,\,\beta=1.0$ & 86.6 & 8.1 & 5.4 \\
$\alpha=0.5,\,\beta=0.5$ & 99.6 & 0.4 & 0.0 \\
$\alpha=2.0,\,\beta=0.5$ & 95.4 & 3.5 & 1.1 \\
$\alpha=0.5,\,\beta=2.0$ & 98.8 & 1.2 & 0.0 \\
\bottomrule
\end{tabular}}
\end{table}


The two weight pairs with the highest illegal rates
($\alpha=1.0,\,\beta=0.0$ and $\alpha=0.0,\,\beta=1.0$) penalise a
single objective most aggressively, which reduces the relative
attractiveness of certain legal actions and encourages broader
exploration early in training.
Despite this, both configurations converge to 100\% acceptance,
confirming that normalisation is sufficient to make legal behaviour
the dominant strategy.

\begin{figure}[H]
    \centering
    \includegraphics[width=0.8\linewidth]{%
        Pictures/Expense_Report_illegal_curves_normalisation.png}
    \caption{Evolution of the illegal-action ratio during training after
             reward normalisation, across all weight configurations.}
    \label{fig:training_normalisation}
\end{figure}

% ------------------------------------------------------------
\subsection{Pareto Recovery Analysis}
\label{sec:pareto_recovery}

Having established that the framework guarantees compliance and that the
agent converges to legal behaviour, this section analyses whether the
recovered policies are not only valid but also optimal, that is,
whether they match the exact Pareto front produced by the CSP+COP solver (see Figure \ref{fig:expense_pareto_scatter}) .

Table~\ref{tab:expense_summary} summarises the acceptance rate, median
cost and duration, and the proportion of episodes converging to each
Pareto point across all six weight configurations.

\begin{table}[H]
\centering
\caption{Training summary after reward normalisation.}
\label{tab:expense_summary}
{\renewcommand{\arraystretch}{1.5}
\setlength{\tabcolsep}{5pt}
\begin{tabular}{ccrrrcc}
\toprule
\textbf{$\alpha$} & \textbf{$\beta$} & \textbf{Accept} & \textbf{Median cost (€)} & \textbf{Median dur.\ (min)}
& \textbf{\% P1} & \textbf{\% P2} \\
\midrule
0.0 & 0.0 & 100\% & 735 & 9{,}900  & 41\% & 1\%  \\
1.0 & 0.0 & 100\% & 735 & 10{,}500 & 44\% & 0\%  \\
0.0 & 1.0 & 100\% & 710 & 2{,}700  & 0\%  & 98\% \\
0.5 & 0.5 & 100\% & 710 & 2{,}700  & 2\%  & 79\% \\
2.0 & 0.5 & 100\% & 710 & 2{,}700  & 2\%  & 95\% \\
0.5 & 2.0 & 100\% & 710 & 2{,}700  & 0\%  & 90\% \\
\bottomrule
\end{tabular}}
\end{table}

The results confirm that the agent recovers both Pareto-optimal traces
across the full range of weight configurations.
The duration focused configuration $(\alpha=0,\,\beta=1)$ converges
almost entirely to P2, the shorter Withdraw path (98\% of episodes).
The cost focused configuration $(\alpha=1,\,\beta=0)$ shows the
strongest preference for P1, the cheaper but longer approval path
(44\% of episodes).
Collectively, both points are reachable through scalarisation, confirming that the \ac{RL} agent approximates the full Pareto front.

One observable pattern is a systematic preference for P2 over P1 even in configurations that are not strongly duration focused.
This bias is examined in Section~\ref{sec:discussion_reward_design} in the context of the per-step penalty structure.

\begin{figure}[H]
    \centering
    \includegraphics[width=0.9\linewidth]{Pictures/expense_pareto_scatter.png}
    \caption{Cost--duration distribution of accepting episodes after
             reward normalisation, across all weight configurations.}
    \label{fig:expense_pareto_scatter}
\end{figure}

% ------------------------------------------------------------
\subsection{Recovered Pareto Traces}
\label{sec:pareto_traces}

To confirm correctness at the trace level,
Table~\ref{tab:pareto_traces} shows the step-by-step execution sequences
for both Pareto-optimal traces as recovered by the \ac{RL} agent.
The sequences match the CSP+COP reference solutions of Diaz et al.\
exactly, validating the framework at the level of individual process
executions.

Notably, both traces are recoverable under the pure duration-minimising
configuration $(\alpha=0.0,\,\beta=1.0)$, indicating that a single
scalarisation is sufficient to cover the entire Pareto front for this
process.

\begin{table}[H]
\centering
\caption{Pareto-optimal trace sequences recovered by the \ac{RL} agent
         (legal steps only).
         Costs in euros (€); durations in minutes.
         $\checkmark$ denotes the accepting state.
         Step reward computed under $\alpha=0.0$, $\beta=1.0$.}
\label{tab:pareto_traces}
{\renewcommand{\arraystretch}{1.5}
\setlength{\tabcolsep}{6pt}
\begin{tabular}{c l r r r}
\toprule
\textbf{Step} & \textbf{Event} & \textbf{Cost (€)} &
\textbf{Duration (min)} & \textbf{Step reward} \\
\midrule
\multicolumn{5}{l}{\textbf{P1}\enspace \texttt{FillOut $\to$ Approve $\to$ PayOut}
  \quad\textit{(total cost: 50, total duration: 10{,}500\,min)}} \\
\midrule
1 & \texttt{FillOutExpenseReport} & 10  & 1{,}500 & $+0.79$ \\
2 & \texttt{Approve}              & 25  & 7{,}200 & $\phantom{+}0.00$ \\
3 & \texttt{PayOut}\;$\checkmark$ & 15  & 1{,}800 & $+100.75$ \\
\midrule
\multicolumn{2}{l}{\textit{Trace total}} &
  \textbf{50} & \textbf{10{,}500} & \\
\midrule
\multicolumn{5}{l}{\textbf{P2}\enspace \texttt{FillOut $\to$ Withdraw}
  \quad\textit{(total cost: 710, total duration: 2{,}700\,min)}} \\
\midrule
1 & \texttt{FillOutExpenseReport}                & 10  & 1{,}500 & $+0.79$ \\
2 & \texttt{WithdrawExpenseReport}\;$\checkmark$ & 700 & 1{,}200 & $+100.83$ \\
\midrule
\multicolumn{2}{l}{\textit{Trace total}} &
  \textbf{710} & \textbf{2{,}700} & \\
\bottomrule
\end{tabular}}
\end{table}

The step rewards in Table~\ref{tab:pareto_traces} follow directly from
the normalised reward of Equation~\ref{eq:reward_total}
(Section~\ref{sec:reward_normalised}); Appendix~\ref{app:expense_step_reward}
works through the computation explicitly for both traces.

% ===================================================

\section{Loan Application: Scalability and Role-Conditioned Trade-offs}
\label{sec:results_loanapp}

This section evaluates the full framework on the role-conditioned Loan
Application benchmark, extending the scale and richness of the
evaluation beyond the five-event Expense Report case: whether the
scalarised reward produces meaningful cost--duration trade-offs in a
larger action space, and what role-conditioned allocation adds once
resource scarcity is introduced. This is the same benchmark
Chapter~\ref{chap:internal_validation} uses to isolate the
contribution of shielding and learning (EQ5), analysed here from a
different angle; the full-framework Pareto front characterised below
is the same one that appears there as the Safe \ac{RL} arm of that
ablation, so compliance behaviour is not re-analysed here. Whether
these findings generalise to process structures never seen during
development is addressed separately on the mined-graph grid of
Chapter~\ref{chap:external_validation}. Omitted diagnostics (reward
convergence, accept rate, shielding activity, clean episodes) are
reported in Appendix~\ref{app:loanapp_diagnostics}.

The process contains 12 events; with the Junior/Expert role profiles,
the action space expands to 24 event--role actions, while the
observation remains the encoded \ac{DCR} marking and role choice
affects only the cost and duration of an executed event.

\paragraph{Two levels of resource modelling.}
The evaluation proceeds in two levels. The first treats Junior and
Expert as always-available role profiles, under which role selection
is separable across events and does not, on its own, constitute
resource allocation. The second introduces scarcity through a per-case
Expert capacity budget, turning role selection into an allocation
problem under opportunity cost. Neither level models inter-case
contention, queues, or utilisation, for the same reason already noted
for this variant in Section~\ref{sec:loan_dataset}; the comparison with
simulation-based approaches such as Kirchdorfer et
al.~\cite{KirchdorferLukasFromPolicies} and
Doumeni~\cite{ArtemisDoumeniParetoAgentSimulator} therefore remains
qualitative, not numerical.

% ------------------------------------------------------------------
\subsection{Pareto Front Recovery (Unlimited Capacity)}
\label{sec:loanapp_pareto}

The Pareto front was extracted from accepting traces generated during
the final 20\,\% of training, with both roles always available. As in
the other Pareto analyses, the unweighted baseline $(\alpha=0,\beta=0)$
is excluded from the dominance computation, since it does not express a
cost--duration preference. The complete list of Pareto-optimal points is
provided in Appendix~\ref{app:loanapp_pareto_points}
(Table~\ref{tab:app_loanapp_pareto_points}).

Figure~\ref{fig:loanapp_pareto_balanced} shows the recovered front. The
agent finds 26 non-dominated cost--duration points, spanning a cost
range from 1\,092 to 3\,120 and a duration range from 408 to 1\,632.
The black stars denote Pareto-optimal accepting traces recovered by the
full framework. The orange diamond marks the most balanced point on the
front, selected as the solution closest to the ideal
minimum-cost/minimum-duration corner after min--max normalising both
objectives. This criterion follows the same interpretation used by
Kirchdorfer et al.~\cite{KirchdorferLukasFromPolicies}: rather than
choosing one extreme, it identifies the compromise point with the
smallest joint distance to the best observed value of each objective.

\textit{This is the single-seed (\(s=1\)) instance of the front. Its
stability across four independent seeds (\(26.25 \pm 0.5\) points, with
near-coincident range and shape) is established in
Section~\ref{sec:policy_generalisation}, and its dependence on both the
shield and the learning component, as opposed to either alone, is
established by the ablation in Section~\ref{sec:shield_contribution}.}

Table~\ref{tab:loanapp_pareto_summary} summarises the three
representative regions of the recovered front: the minimum-cost
solution, the minimum-duration solution, and the balanced compromise.
The full 26-point table is reported separately in
Appendix~\ref{app:loanapp_pareto_points} to keep the main result
focused on the structure of the front rather than on the enumeration of
all points.

\begin{figure}[htpb]
  \centering
  \begin{minipage}{0.50\textwidth}
    \centering
    \includegraphics[width=\textwidth]{Pictures/LoanApp_pareto_balanced.png}
    \par\footnotesize(a) Full recovered front.
  \end{minipage}%
  \hfill
  \begin{minipage}{0.50\textwidth}
    \centering
    \includegraphics[width=\textwidth]{Pictures/LoanApp_pareto_balanced_zoom.png}
    \par\footnotesize(b) Zoomed-in central cluster.
  \end{minipage}
  \caption{Pareto front recovered by the full \ac{PPO}+\ac{DCR}
  framework on the role-conditioned Loan Application benchmark, with
  both roles always available. Points correspond to accepting traces
  from the final 20\,\% of training, with the unweighted baseline
  excluded from dominance computation. Black stars denote non-dominated
  traces; the orange diamond marks the most balanced point, selected by
  minimum Euclidean distance to the ideal point after min--max
  normalisation. (b) zooms into the central cluster where most
  intermediate Pareto points concentrate.}
  \label{fig:loanapp_pareto_balanced}
\end{figure}



\begin{table}[H]
\centering
\caption{Summary of representative Pareto points recovered on the
role-conditioned Loan Application benchmark (unlimited capacity). The
balanced point is selected as the point closest to the ideal
minimum-cost/minimum-duration corner after min--max normalisation.}
\label{tab:loanapp_pareto_summary}
{\renewcommand{\arraystretch}{1.5}
\setlength{\tabcolsep}{8pt}
\begin{tabular}{lrrl}
\toprule
\textbf{Point} & \textbf{Cost (a.u.)} & \textbf{Duration (min)} & \textbf{Configuration} \\
\midrule
Minimum cost      & 1\,092.0 & 1\,632.0 & $\alpha=1.0,\ \beta=0.0$ \\
Balanced          & 2\,044.9 & 1\,036.5 & $\alpha=0.5,\ \beta=0.5$ \\
Minimum duration  & 3\,120.0 &   408.0 & $\alpha=0.0,\ \beta=1.0$ \\
\bottomrule
\end{tabular}}
\end{table}

To check whether this extreme of the front is a true optimum rather
than a training artefact, the minimum-cost point was verified
independently against the exact CSP+COP solver. Role choice was
removed for this check, reducing every event to a single, flat cost
and duration value; under this simplification, the solver's Pareto
front over the same five events (\textit{Check application form
completeness}, \textit{Appraise property}, \textit{Check credit
history}, \textit{AML check}, \textit{Assess loan risk}) contains
exactly one non-dominated point, cost\,=\,1\,560, duration\,=\,816,
reached in 5.6\,s and proven optimal. A single point is the expected
outcome here, not a weakness of the check: unlike Expense Report,
this flat sub-problem admits no genuine cost--duration trade-off.

Rescaling both figures by the Junior role multiplier ($\times0.7$
cost, $\times2.0$ duration) gives exactly $(1\,092.0,\,1\,632.0)$,
matching the recovered minimum-cost point on the 26-point front
exactly. This confirms that this extreme of the front is the true
global optimum, not an artefact of training.

This check is nonetheless narrow in scope: it covers only a flat,
role-free problem with a single-point front. Unlike Expense Report,
where both Pareto-optimal points are independently confirmed by the
solver (Section~\ref{sec:expense_baseline}), no exact reference exists
for the remaining 25 points of the Loan Application front, including
its minimum-duration extreme and every intermediate point. ``Exact
recovery'' in this thesis is therefore verified at the level of a
full Pareto front only on benchmarks with small, degenerate fronts (at
most two points, or one in this flat sub-problem); on the only
non-trivial front recovered, only one extreme is anchored to a proven
optimum.


The location of these three points confirms that the scalarisation
weights steer the trained policies towards different regions of the
trade-off surface. The cost-focused configuration reaches the cheapest
valid executions, the duration-focused configuration reaches the
fastest valid executions, and the balanced configuration contributes a
central compromise. A notable pattern, visible in the complete table in
Appendix~\ref{app:loanapp_pareto_points}, is that 15 of the 26 Pareto
points are produced by the balanced configuration
$(\alpha=0.5,\beta=0.5)$ rather than by the two extreme configurations.
This suggests that balanced rewards maintain broader exploration of the
valid trace space, whereas specialised objectives converge more quickly
towards one end of the front. This behaviour is revisited in
Section~\ref{sec:discussion_reward}.

% ------------------------------------------------------------------
\subsection{Effect of Objective Weights}
\label{sec:loanapp_weight_effects}

Table~\ref{tab:loanapp_relative} reports the median cost and duration
of accepting traces in the converged phase for each weight
configuration, expressed relative to the unweighted baseline;
additional training diagnostics for the same runs are reported in
Appendix~\ref{app:loanapp_diagnostics}. The results show the expected
directional behaviour: duration-dominant configurations reduce
duration by 37.7\,\%, at a cost increase of 22--27\,\%, while
cost-dominant configurations reduce cost by 32.4\,\% at the price of
longer duration. This unconstrained effect is superseded by the
richer, heuristic-compared analysis under scarcity in
Section~\ref{sec:loanapp_capacity} below.

\begin{table}[H]
\centering
\caption{Median cost and duration in the converged phase (last 20\,\%
of training), with relative change against the unweighted baseline
$(\alpha=0,\beta=0)$.}
\label{tab:loanapp_relative}
{\renewcommand{\arraystretch}{1.2}
\setlength{\tabcolsep}{4pt}
\begin{tabular}{lrrrr}
\toprule
\textbf{Configuration} & \textbf{Median cost (a.u.)} & \textbf{Median duration (min)}
  & \textbf{$\Delta$ cost} & \textbf{$\Delta$ duration} \\
\midrule
$\alpha=0.0,\;\beta=0.0$ & 2\,248.1 & 1\,087.5 & --       & --       \\
$\alpha=0.0,\;\beta=1.0$ & 2\,748.6 &   678.0 & +22.3\,\% & -37.7\,\% \\
$\alpha=0.5,\;\beta=0.5$ & 1\,949.1 & 1\,149.0 & -13.3\,\% & +5.7\,\%  \\
$\alpha=0.5,\;\beta=2.0$ & 2\,858.5 &   678.0 & +27.2\,\% & -37.7\,\% \\
$\alpha=1.0,\;\beta=0.0$ & 1\,520.1 & 1\,513.5 & -32.4\,\% & +39.2\,\% \\
$\alpha=2.0,\;\beta=0.5$ & 1\,520.1 & 1\,399.5 & -32.4\,\% & +28.7\,\% \\
\bottomrule
\end{tabular}}
\end{table}

% ------------------------------------------------------------------
\subsection{Role Specialisation Patterns}
\label{sec:loanapp_roles}

Beyond the aggregate cost--duration trade-off, the role-conditioned
action space lets the policy express \emph{which} events it staffs with
each role, and how that choice responds to the objective weights.
Figure~\ref{fig:loanapp_role_slopes} reports the Expert-assignment rate
per event under the cost-focused $(\alpha{=}1,\beta{=}0)$ and
duration-focused $(\alpha{=}0,\beta{=}1)$ regimes; the slope of each
line measures how strongly an event's staffing responds to the
objective.

\begin{figure}[H]
  \centering
  \includegraphics[width=0.8\textwidth]{Pictures/LoanApp_role_slopegraph.png}
  \caption{Expert-assignment rate per event under the cost-focused and
  duration-focused regimes. Each line is one event; a steep slope
  indicates strong sensitivity of the staffing decision to the
  objective, a flat line indicates a role-indifferent activity.}
  \label{fig:loanapp_role_slopes}
\end{figure}

Three qualitative categories emerge. High-sensitivity
events such as \emph{AML check} and \emph{Appraise property} swing
sharply between Junior and Expert dominance as the objective changes.
Lower-impact events respond weakly. \emph{Assess loan risk} is a
boundary case, remaining close to a balanced 50\,\% Expert assignment
irrespective of the weights.

The mechanism behind the three categories, and its implications for
process management, are analysed in Section~\ref{sec:discussion_roles}.
Their stability across random seeds, at both the front level and the
per-event role-assignment level, is established in
Section~\ref{sec:policy_generalisation}.

% ------------------------------------------------------------------
\subsection{Capacity-constrained Role Allocation}
\label{sec:loanapp_capacity}

The analysis so far assumes both roles are always available, under
which role selection is separable across events
(Section~\ref{sec:results_loanapp}). This section removes that
assumption by limiting the Expert role to a
per-case budget of $k$ assignments, as defined in
Section~\ref{sec:loan_dataset}. A finite budget couples the decisions,
so the policy must allocate a scarce resource rather than apply a
per-event rule.

The constraint and the optimisation are kept separate. The budget limit
is enforced by the same shield that enforces \ac{DCR} enablement: an
Expert action with no remaining budget is intercepted without modifying
the marking, so the compliance guarantee of contribution~C1 extends
unchanged to the capacity constraint. The shield, however, only blocks;
it does not decide where the budget is spent. That decision is learned.
The remaining budget is part of the observation, so the policy can
condition on it, and the reward signal teaches the agent that spending
the scarce Expert early on a low-impact event leaves none for a
high-impact one later. Allocating the budget well is therefore a
sequential credit-assignment problem of the same kind the agent solves
when recovering the Pareto front, and one that a myopic per-event
heuristic cannot address by construction.

\paragraph{Relation to Kirchdorfer et al.'s evaluation design.}
Because the Loan Application costs are synthetic and the metrics and
granularity differ from prior work, the cost and duration figures
recovered here are not numerically comparable to those of Kirchdorfer et
al.~\cite{KirchdorferLukasFromPolicies} or
Doumeni~\cite{ArtemisDoumeniParetoAgentSimulator}; that comparison
remains qualitative and is developed in Chapter~\ref{chap:discussion}.
This experiment is best understood as a methodological replication
rather than as external validation in the sense of
Chapter~\ref{chap:external_validation}: it reproduces the evaluation
design of prior resource-allocation work, in which an optimised policy
is benchmarked against standard allocation heuristics, and applies it
independently on a \ac{DCR}-shielded learning agent, on the same
internal benchmark used throughout this chapter and
Chapter~\ref{chap:internal_validation}. The heuristics used here,
Always-Junior, Always-Expert, and greedy-per-event, are the same family
of cheapest, fastest, and myopic-greedy references that Kirchdorfer et
al.\ benchmark against. The question is therefore not whether the
recovered numbers match theirs, but whether the same qualitative
conclusion holds, namely that a learned allocation policy improves on
naive heuristics under scarcity.

The experiment is organised around three questions:
\begin{description}
  \item[Q1.] How does scarcity reshape the recovered front relative to
  the unconstrained case of Section~\ref{sec:loanapp_pareto}?
  \item[Q2.] When Expert capacity is scarce, does the learned policy
  concentrate it on the high-sensitivity events identified in
  Section~\ref{sec:loanapp_roles}?
  \item[Q3.] Does the learned policy improve on the standard allocation
  heuristics, as an independent replication of the
  optimised-policy-versus-heuristics comparison reported by Kirchdorfer
  et al.?
\end{description}

\paragraph{Effect of the budget on the front (Q1).}
Table~\ref{tab:q3_heuristics} reports the recovered front for each
budget $k \in \{1,2,3,\infty\}$, alongside the three allocation
heuristics evaluated at the same budgets (the full learned-vs-heuristic
comparison is developed under Q3 below; this paragraph focuses on the
\textbf{Learned front} column). The $k=\infty$ row reproduces the
unconstrained front of Section~\ref{sec:loanapp_pareto}. As the budget
tightens, the front both shrinks (fewer Pareto points: 26 at
$k=\infty$ down to 6 at $k=1$) and shifts toward higher minimum
duration (408 at $k=\infty$ versus 1\,245 at $k=1$), while the
minimum-cost extreme (1\,092) is unaffected at every budget level, since
the cheapest trace never needs Expert regardless of how much capacity is
available. Always-Junior and Always-Expert are summarised by their
accept rate and median outcome rather than a range, since neither uses
$(\alpha,\beta)$ to target a trade-off and each instead produces a wide
distribution of episode outcomes under random event ordering.

The same contraction is visible at the level of the full cost and
duration distributions, not just the front's extremes.
Figure~\ref{fig:loanapp_capacity_box} repeats the weight-pair comparison
of Section~\ref{sec:loanapp_weight_effects} under each budget
$k \in \{1,2,3\}$. At $k=1$ the six weight pairs collapse into nearly
indistinguishable cost and duration distributions: with only one Expert
assignment available per case, the weights have almost nothing left to
act on. As $k$ grows to 2 and 3, the distributions separate cleanly by
objective, cost-dominant pairs shifting to visibly lower cost,
duration-dominant pairs to visibly lower duration, recovering the same
directional pattern as the unconstrained case. Scarcity, in other
words, does not just shrink the front, it narrows how much the
scalarisation weights can differentiate outcomes at all.

\begin{figure}[htpb]
  \centering
  \begin{minipage}{0.85\textwidth}
    \centering
    \includegraphics[width=\textwidth]{Pictures/loanapp_capacity_cost_boxplot.png}
    \par\footnotesize(a) Total cost.
  \end{minipage}

  \medskip
  \begin{minipage}{0.85\textwidth}
    \centering
    \includegraphics[width=\textwidth]{Pictures/loanapp_capacity_dur_boxplot.png}
    \par\footnotesize(b) Total duration.
  \end{minipage}
  \caption{Distribution of total (a) cost and (b) duration across
  accepting traces in the final 20\,\% of training, grouped by
  scalarisation weight pair, under each Expert capacity budget
  $k \in \{1,2,3\}$. At $k=1$ the weight pairs are barely
  distinguishable; the separation by objective sharpens as the budget
  grows.}
  \label{fig:loanapp_capacity_box}
\end{figure}


\begin{table}[htpb]
\centering
\footnotesize
\caption{Learned policy versus allocation heuristics at each Expert
budget $k$, given as (Pareto points, cost range / duration range) for
the learned and greedy fronts, and median (cost, duration) for
Always-Expert; cost in a.u., duration in min throughout. All heuristic
and learned configurations reach 100\,\% acceptance at every budget.
Always-Junior never selects Expert and is therefore unaffected by $k$:
median (cost, duration) $=(3\,126,\,4\,683)$ throughout.}
\label{tab:q3_heuristics}
{\renewcommand{\arraystretch}{1.5}
\setlength{\tabcolsep}{5pt}
\begin{tabular}{lrrr}
\toprule
\textbf{Budget} & \textbf{Learned front} & \textbf{Greedy front} &
\textbf{Always-Expert} \\
\midrule
$k=1$      & 6 pts, 1\,092--1\,793 / 1\,245--1\,632 & 2 pts, 1\,092--1\,191 / 1\,550--1\,632 & median (3\,248, 4\,566) \\
$k=2$      & 14 pts, 1\,092--2\,454 / 879--1\,632 & 5 pts, 1\,092--1\,892 / 1\,162--1\,632 & median (3\,541, 4\,380) \\
$k=3$      & 21 pts, 1\,092--2\,868 / 616--1\,632 & 7 pts, 1\,092--2\,553 / 796--1\,632 & median (3\,996, 4\,200) \\
$k=\infty$ & 26 pts, 1\,092--3\,120 / 408--1\,632 & 5 pts, 1\,228--3\,120 / 408--1\,910 & median (8\,932, 1\,171) \\
\bottomrule
\end{tabular}}
\end{table}

Figure~\ref{fig:loanapp_capacity_pareto} visualises this contraction
directly: the $k=\infty$ curve is the same 26-point front shown in
Figure~\ref{fig:loanapp_pareto_balanced}, and each tighter budget traces
a strict sub-front of it, pivoting around the shared minimum-cost
extreme while progressively losing the lower-duration, higher-cost
region.

\begin{figure}[htpb]
  \centering
  \includegraphics[width=0.75\textwidth]{Pictures/LoanApp_capacity_pareto_combined.png}
  \caption{Recovered Pareto front under each Expert capacity budget $k$.
  The $k=\infty$ curve reproduces the unconstrained front of
  Figure~\ref{fig:loanapp_pareto_balanced}; tighter budgets trace
  progressively shorter sub-fronts, anchored at the shared minimum-cost
  point.}
  \label{fig:loanapp_capacity_pareto}
\end{figure}

\paragraph{Where the scarce Expert is spent (Q2).}
Figure~\ref{fig:loanapp_capacity_allocation} shows, for each budget, the
share of the Expert allocation directed to each event, with the
high-sensitivity events of Section~\ref{sec:loanapp_roles} marked. Three
distinct allocation patterns emerge as $k$ varies. First, a monotonic
group (AML check, Appraise property) receives a steadily growing share
of Expert as the budget grows (9\%$\to$35\%$\to$41\%$\to$48\% for AML
check), consistent with these being the high-sensitivity events already
identified in Section~\ref{sec:loanapp_roles}: the policy reinforces
them first whenever capacity allows. Second, an indifferent event
(Cancel application) holds at approximately 50\% Expert assignment
regardless of $k$, including under the most extreme scarcity ---
stronger than the boundary behaviour of Assess loan risk, which is
indifferent only when capacity is unlimited and collapses to 6\% under
$k=1$. Third, a front-loaded event (Design loan offer) shows the
opposite of the monotonic pattern: under extreme scarcity ($k=1$) it
absorbs a disproportionate 70\% of the available budget, a priority that
dilutes as capacity grows (70\%$\to$61\%$\to$52\%$\to$41\%). This third
pattern was not anticipated and indicates that scarcity does not simply
scale down the unconstrained allocation uniformly --- it changes which
events the policy treats as priorities.

\begin{figure}[htpb]
  \centering
  \includegraphics[width=\textwidth]{Pictures/loanapp_capacity_allocation.png}
  \caption{Distribution of the scarce Expert budget across events under
  each capacity limit $k$. The high-sensitivity events identified in
  Section~\ref{sec:loanapp_roles} are highlighted.}
  \label{fig:loanapp_capacity_allocation}
\end{figure}

The pooled allocation pattern of Figure~\ref{fig:loanapp_capacity_allocation}
aggregates over many episodes; Appendix~\ref{app:loanapp_capacity_trace}
shows the same mechanism at the level of a single representative trace
per budget, confirming that the agent allocates the scarce Expert
budget selectively where it most affects the trade-off, rather than
exhausting it as a fixed quota.

\paragraph{Learned policy versus allocation heuristics (Q3).}
The comparison against the three heuristics is the methodological core
of this replication: it reproduces, on an independent method and
benchmark, the optimised-policy-versus-heuristics design of Kirchdorfer
et al.~\cite{KirchdorferLukasFromPolicies}. Table~\ref{tab:q3_heuristics}
(introduced under Q1 above) compares the methods at every budget $k$.


Three findings follow. First, the learned policy extends the recovered
front further than the greedy heuristic at every budget: the two track
each other closely in the region where both operate
(Figure~\ref{fig:q3_heuristics}b), but the learned policy alone
reaches the cheapest and the most duration-efficient extremes, recovering
the qualitative conclusion of Kirchdorfer et al.\ from a different
paradigm while also retaining the compliance guarantee of
Section~\ref{sec:shield_contribution}, which neither heuristic provides
by construction. Second, both non-weight-conditioned heuristics are far
worse than either optimised approach at every budget (median duration
above 4\,000 against a learned-policy maximum of 1\,632), confirming
that ignoring cost and duration entirely is costly even when
feasibility is not at stake. Third, and the most informative
comparison of the three: Always-Expert is barely better than
Always-Junior, despite having access to a resource Always-Junior
never uses at all. At $k=2$, for example, Always-Expert pays $415$
more in median cost ($3\,541$ against $3\,126$) to save only $303$ in
median duration ($4\,380$ against $4\,683$). The learned policy and
Greedy get a far better return from the same budget: at $k=2$, the
learned front alone reaches a duration as low as $879$ at a cost of
only $2\,454$ (Table~\ref{tab:q3_heuristics}), well below
Always-Expert's duration at comparable or lower cost.

The difference is not whether the Expert budget is available, but
\emph{where} it is spent. Always-Expert assigns it to whichever event
happens to be selected first, with no regard for which events the
process benefits from most, and captures almost none of its value as
a result. A fixed rule with no notion of cost or duration cannot make
that distinction; the learned policy and Greedy can.

\begin{figure}[htpb]
  \centering
  \begin{minipage}{0.47\textwidth}
    \centering
    \includegraphics[width=\textwidth]{Pictures/q3_heuristics_k1.png}
    \par\footnotesize(a) $k=1$.
  \end{minipage}\hfill
  \begin{minipage}{0.47\textwidth}
    \centering
    \includegraphics[width=\textwidth]{Pictures/q3_heuristics_k2.png}
    \par\footnotesize(b) $k=2$.
  \end{minipage}

  \medskip
  \begin{minipage}{0.47\textwidth}
    \centering
    \includegraphics[width=\textwidth]{Pictures/q3_heuristics_k3.png}
    \par\footnotesize(c) $k=3$.
  \end{minipage}
  \caption{Learned policy versus role-allocation heuristics at budgets
  (a)~$k=1$, (b)~$k=2$, and (c)~$k=3$. At every budget, the learned
  policy (blue) extends beyond the greedy heuristic (green) at both
  ends of the front, while Always-Expert (red) and Always-Junior
  (purple), each shown at their median outcome, remain far from
  either.}
  \label{fig:q3_heuristics}
\end{figure}
